\begin{examplebox}
\textbf{例1（必考陷阱）：竖直切线不能遗漏。}

过 $M(2,2)$ 作椭圆 $\frac{x^2}{4}+y^2=1$ 的切线，求所有切线。

\textit{解：} $a^2=4,b^2=1$，代入 (C009-1)：
\[
(4-4)k^2+8k+(1-4)=0\quad\Longrightarrow\quad k=\frac38.
\]
有限斜率的切线为 $y-2=\frac38(x-2)$，即
\[
\boxed{y=\frac38x+\frac54}.
\]
但 $x_0=a=2$，还应补上长轴顶点处的竖直切线
\[
\boxed{x=2}.
\]
\textbf{答：}共两条；“两个切线斜率的乘积”在本题中\textbf{无定义}。

\medskip
\textbf{例2（韦达定理）：两条有限斜率的和与积。}

过 $M(0,2)$ 作椭圆 $\frac{x^2}{4}+y^2=1$ 的切线。
\[
4k^2-3=0\quad\Longrightarrow\quad
\boxed{k_1=\frac{\sqrt3}{2},\quad k_2=-\frac{\sqrt3}{2}}.
\]
故 $k_1+k_2=0,\ k_1k_2=-\frac34$；两条切线分别为
\[
\boxed{y=\frac{\sqrt3}{2}x+2,\qquad y=-\frac{\sqrt3}{2}x+2}.
\]
\textit{数形验证：}两切点为 $(-\sqrt3,\frac12)$ 与 $(\sqrt3,\frac12)$；切点弦是 $y=\frac12$。

\medskip
\textbf{例3（椭圆上的点）：一个重根对应一条切线。}

在 $P(1,\frac{\sqrt3}{2})$ 处作椭圆 $\frac{x^2}{4}+y^2=1$ 的切线。

先检验 $\frac14+\frac34=1$，故 $P$ 在椭圆上。
切线斜率
\[
\boxed{k=-\frac{b^2x_0}{a^2y_0}=-\frac{\sqrt3}{6}}.
\]
切线方程为
\[
\boxed{\frac{x}{4}+\frac{\sqrt3}{2}y=1}.
\]
代入斜率方程也只得到一个重根，不代表两条不同切线。

\medskip
\textbf{例4（综合）：两切线垂直，求切线交点轨迹。}

设过 $M(x_0,y_0)$ 引椭圆 $\frac{x^2}{a^2}+\frac{y^2}{b^2}=1$ 的两条切线互相垂直。

\textit{情况一：} $x_0^2\ne a^2$，两条切线都不竖直，故 $k_1k_2=-1$。利用韦达定理，
\[
\frac{b^2-y_0^2}{a^2-x_0^2}=-1
\quad\Longrightarrow\quad\boxed{x_0^2+y_0^2=a^2+b^2}.
\]
\textit{情况二：} $x_0=\pm a$。一条切线为 $x=x_0$；要垂直，另一条必须水平，即 $k=0$。
由 $2x_0y_0k+b^2-y_0^2=0$ 得 $y_0=\pm b$，仍在上述圆上。

\textit{反向验证：}圆 $x_0^2+y_0^2=a^2+b^2$ 上每一点都在椭圆外；当 $x_0^2\ne a^2$ 时，两有限斜率乘积为 $-1$，而 $x_0=\pm a$ 时为一横一直。
因此轨迹恰为
\[
\boxed{x_0^2+y_0^2=a^2+b^2}\qquad\text{（蒙日圆）}.
\]
例如 $\frac{x^2}{4}+y^2=1$ 时，$M(2,1)$ 的两条切线分别是 $x=2$ 和 $y=1$，确实互相垂直。

\medskip
\textbf{例5（全集性质）：指定斜率，反求全部切线。}

求椭圆 $\frac{x^2}{4}+y^2=1$ 中斜率为 $2$ 的所有切线。

由 $m^2=a^2k^2+b^2=4\times 4+1=17$，故
\[
\boxed{y=2x+\sqrt{17},\qquad y=2x-\sqrt{17}}.
\]
可见椭圆并没有限制\textbf{有限切线斜率}的取值范围：$k$ 可取一切实数。
\end{examplebox}
